% ========================================================
% BÀI TẬP TOÁN 5 - KẾT NỐI TRI THỨC VỚI CUỘC SỐNG
% Bài 71: Ôn tập hình học
% Định dạng: BTTeX Part III - Trắc nghiệm điền số ngắn
% ========================================================

% Câu 1
\begin{ex}
  Một hình tam giác có độ dài đáy $12\text{ cm}$ và chiều cao tương ứng $8\text{ cm}$. Diện tích tam giác là bao nhiêu xăng-ti-mét vuông?
  \shortans{48}
  \loigiai{
    $(12 \times 8) : 2 = 48\text{ cm}^2$.
  }
\end{ex}

% Câu 2
\begin{ex}
  Một hình thang có đáy lớn $15\text{ cm}$, đáy bé $9\text{ cm}$, chiều cao $6\text{ cm}$. Diện tích hình thang là bao nhiêu xăng-ti-mét vuông?
  \shortans{72}
  \loigiai{
    $(15 + 9) \times 6 : 2 = 72\text{ cm}^2$.
  }
\end{ex}

% Câu 3
\begin{ex}
  Một hình tròn có bán kính $10\text{ cm}$. Chu vi của hình tròn là bao nhiêu xăng-ti-mét khi lấy $\pi = 3{,}14$ và nhân với 10?
  \shortans{628}
  \loigiai{
    Chu vi: $2 \times 10 \times 3{,}14 = 62{,}8\text{ cm}$. Nhân 10 được 628.
  }
\end{ex}

% Câu 4
\begin{ex}
  Một hình tròn có bán kính $10\text{ cm}$. Diện tích của hình tròn là bao nhiêu xăng-ti-mét vuông (lấy $\pi = 3{,}14$)?
  \shortans{314}
  \loigiai{
    $10 \times 10 \times 3{,}14 = 314\text{ cm}^2$.
  }
\end{ex}

% Câu 5
\begin{ex}
  Một hình lập phương có cạnh $5\text{ cm}$. Diện tích xung quanh của hình lập phương là bao nhiêu xăng-ti-mét vuông?
  \shortans{100}
  \loigiai{
    $5 \times 5 \times 4 = 100\text{ cm}^2$.
  }
\end{ex}

% Câu 6
\begin{ex}
  Ở câu trên, diện tích toàn phần của hình lập phương là bao nhiêu xăng-ti-mét vuông?
  \shortans{150}
  \loigiai{
    $5 \times 5 \times 6 = 150\text{ cm}^2$.
  }
\end{ex}

% Câu 7
\begin{ex}
  Ở câu trên, thể tích của hình lập phương là bao nhiêu xăng-ti-mét khối?
  \shortans{125}
  \loigiai{
    $5 \times 5 \times 5 = 125\text{ cm}^3$.
  }
\end{ex}

% Câu 8
\begin{ex}
  Một hình hộp chữ nhật có chiều dài $8\text{ cm}$, chiều rộng $5\text{ cm}$, chiều cao $4\text{ cm}$. Chu vi mặt đáy là bao nhiêu xăng-ti-mét?
  \shortans{26}
  \loigiai{
    $(8 + 5) \times 2 = 26\text{ cm}$.
  }
\end{ex}

% Câu 9
\begin{ex}
  Ở hình hộp chữ nhật trên, diện tích xung quanh là bao nhiêu xăng-ti-mét vuông?
  \shortans{104}
  \loigiai{
    $26 \times 4 = 104\text{ cm}^2$.
  }
\end{ex}

% Câu 10
\begin{ex}
  Ở hình hộp chữ nhật trên, thể tích là bao nhiêu xăng-ti-mét khối?
  \shortans{160}
  \loigiai{
    $8 \times 5 \times 4 = 160\text{ cm}^3$.
  }
\end{ex}

% Câu 11
\begin{ex}
  Một mảnh vườn hình chữ nhật có chiều dài $25\text{ m}$, chiều rộng $12\text{ m}$. Diện tích mảnh vườn là bao nhiêu mét vuông?
  \shortans{300}
  \loigiai{
    $25 \times 12 = 300\text{ m}^2$.
  }
\end{ex}

% Câu 12
\begin{ex}
  Ở mảnh vườn trên, chu vi là bao nhiêu mét?
  \shortans{74}
  \loigiai{
    $(25 + 12) \times 2 = 74\text{ m}$.
  }
\end{ex}

% Câu 13
\begin{ex}
  Một hình tam giác vuông có hai cạnh góc vuông là $6\text{ cm}$ và $8\text{ cm}$. Diện tích của hình tam giác đó là bao nhiêu xăng-ti-mét vuông?
  \shortans{24}
  \loigiai{
    $(6 \times 8) : 2 = 24\text{ cm}^2$.
  }
\end{ex}

% Câu 14
\begin{ex}
  Một hình thoi có độ dài hai đường chéo là $14\text{ cm}$ và $10\text{ cm}$. Diện tích hình thoi là bao nhiêu xăng-ti-mét vuông?
  \shortans{70}
  \loigiai{
    $(14 \times 10) : 2 = 70\text{ cm}^2$.
  }
\end{ex}

% Câu 15
\begin{ex}
  Một hình tròn có đường kính $20\text{ cm}$. Bán kính của hình tròn là bao nhiêu xăng-ti-mét?
  \shortans{10}
  \loigiai{
    $20 : 2 = 10\text{ cm}$.
  }
\end{ex}

% Câu 16
\begin{ex}
  Một hình thang có diện tích $60\text{ cm}^2$, chiều cao $5\text{ cm}$. Tổng độ dài hai đáy là bao nhiêu xăng-ti-mét?
  \shortans{24}
  \loigiai{
    $(60 \times 2) : 5 = 24\text{ cm}$.
  }
\end{ex}

% Câu 17
\begin{ex}
  Một căn phòng hình hộp chữ nhật dài $6\text{ m}$, rộng $4\text{ m}$, cao $3\text{ m}$. Diện tích xung quanh phòng là bao nhiêu mét vuông?
  \shortans{60}
  \loigiai{
    $(6 + 4) \times 2 \times 3 = 60\text{ m}^2$.
  }
\end{ex}

% Câu 18
\begin{ex}
  Ở căn phòng trên, diện tích trần nhà là bao nhiêu mét vuông?
  \shortans{24}
  \loigiai{
    $6 \times 4 = 24\text{ m}^2$.
  }
\end{ex}

% Câu 19
\begin{ex}
  Một khối gỗ hình lập phương có cạnh $0{,}4\text{ m}$. Thể tích của khối gỗ đó khi nhân với 1\,000 là bao nhiêu?
  \shortans{64}
  \loigiai{
    $0{,}4 \times 0{,}4 \times 0{,}4 = 0{,}064\text{ m}^3 = 64\text{ dm}^3$.
  }
\end{ex}

% Câu 20
\begin{ex}
  Một bể cá hình hộp chữ nhật dài $80\text{ cm}$, rộng $50\text{ cm}$, cao $40\text{ cm}$. Thể tích của bể là bao nhiêu lít nước ($1\text{ lít} = 1\text{ dm}^3$)?
  \shortans{160}
  \loigiai{
    $8 \times 5 \times 4 = 160\text{ dm}^3 = 160$ lít.
  }
\end{ex}

% Câu 21
\begin{ex}
  Một hình tam giác có diện tích $30\text{ cm}^2$, đáy dài $10\text{ cm}$. Chiều cao tương ứng là bao nhiêu xăng-ti-mét?
  \shortans{6}
  \loigiai{
    $(30 \times 2) : 10 = 6\text{ cm}$.
  }
\end{ex}

% Câu 22
\begin{ex}
  Một hình lập phương có diện tích toàn phần là $96\text{ cm}^2$. Diện tích một mặt của hình lập phương là bao nhiêu xăng-ti-mét vuông?
  \shortans{16}
  \loigiai{
    $96 : 6 = 16\text{ cm}^2$.
  }
\end{ex}

% Câu 23
\begin{ex}
  Ở hình lập phương trên, độ dài cạnh là bao nhiêu xăng-ti-mét?
  \shortans{4}
  \loigiai{
    Vì $4 \times 4 = 16$ nên cạnh bằng 4 cm.
  }
\end{ex}

% Câu 24
\begin{ex}
  Ở hình lập phương trên, thể tích là bao nhiêu xăng-ti-mét khối?
  \shortans{64}
  \loigiai{
    $4 \times 4 \times 4 = 64\text{ cm}^3$.
  }
\end{ex}

% Câu 25
\begin{ex}
  Một sân cỏ hình chữ nhật dài $100\text{ m}$, rộng $60\text{ m}$. Xung quanh xây tường rào, để một cổng rộng $4\text{ m}$. Chiều dài tường rào là bao nhiêu mét?
  \shortans{316}
  \loigiai{
    Chu vi: $(100 + 60) \times 2 = 320\text{ m}$. Chiều dài tường rào: $320 - 4 = 316\text{ m}$.
  }
\end{ex}
